% --- Archon physics formalization source begin ---
% archon:physics
% archon:covers PhyXMiniProblems/problem_phyx_mini_0774.lean
% archon:source-report reports/phyx_mini/problem_phyx_mini_0774.source.json

\chapter{Physics problem phyx\_mini\_0774}
\label{ch:PhyXMiniProblems_problem_phyx_mini_0774}

\paragraph{Problem source.}
\#\# Physical scenario

The pulley in figure has radius \textbackslash{}(0.160\textbackslash{}text\{ m\}\textbackslash{}) and moment of inertia \textbackslash{}(0.380\textbackslash{}text\{ kg\}\textbackslash{}cdot\textbackslash{}text\{m\}\textasciicircum{}2\textbackslash{}). The rope does not slip on the pulley rim.

\#\# Question

calculate the speed of the \textbackslash{}(4.00\textbackslash{}text\{ kg\}\textbackslash{}) block just before it strikes the floor.

\#\# Answer choices

- A: A: 3.56m/s

- B: B: 3.07m/s

- C: C: 4.07m/s

- D: D: 4.56m/s

\#\# Figure caption (auxiliary; use the image as primary evidence)

The image shows a simple pulley system. It includes the following elements:

1. **Pulley**: Positioned at the top, fixed to a support or ceiling.
2. **Rope**: Draped over the pulley, connecting two blocks.
3. **Blocks**:
   - **Left block**: Labeled as 4.00 kg, suspended in the air.
   - **Right block**: Labeled as 2.00 kg, resting on the ground.
4. **Distance**: There is a vertical arrow indicating a distance of 5.00 m between the two blocks.
5. **Arrows**: Indicate the direction of motion or forces acting on the blocks; the arrow points downward from the 4.00 kg block suggesting it will move down. 

The setup suggests a scenario likely used to illustrate principles of mechanics, such as tension, acceleration, or gravitational force.

\#\# Dataset metadata

Dynamics; Physical Model Grounding Reasoning, Multi-Formula Reasoning

\paragraph{Recorded answer/context.}
B: B: 3.07m/s

\paragraph{Figure/image path.}
/root/proposal\_for\_physic/science-mango/phyx\_mini\_run/phyx\_data/test\_image/774.png

\paragraph{Formalization target.}
create a compiling Lean file with sorry bodies at `PhyXMiniProblems/problem\_phyx\_mini\_0774.lean`. The Lean declarations must preserve the physical quantities, dimensions or dimensional roles, figure labels, governing-law hypotheses, and final relation expressed by this problem.
Use Mathlib/Physlib names found through LeanExplore where available. If a physics API is missing, introduce faithful local abstractions rather than scalar placeholder aliases.

\begin{theorem}[Physics formalization target]
\label{thm:physics:phyx_mini_0774:target}
\lean{PhyXMiniProblems.ProblemPhyXMini0774.problem_phyx_mini_0774}
\uses{def:physics:phyx-mini-0774:phyxminiproblems-problemphyxmini0774-speedinmeterspersecond, def:physics:phyx-mini-0774:phyxminiproblems-problemphyxmini0774-massivepulleyatwoodsetup, def:physics:phyx-mini-0774:phyxminiproblems-problemphyxmini0774-matchesprimaryfigure, def:physics:phyx-mini-0774:phyxminiproblems-problemphyxmini0774-matchesproblemstatement, def:physics:phyx-mini-0774:phyxminiproblems-problemphyxmini0774-usesidealtextbookreleasemodel, def:physics:phyx-mini-0774:phyxminiproblems-problemphyxmini0774-usesstandardnearearthgravity, def:physics:phyx-mini-0774:phyxminiproblems-problemphyxmini0774-hasphysicalparameters, def:physics:phyx-mini-0774:phyxminiproblems-problemphyxmini0774-satisfiesmassivepulleyenergylaws, def:physics:phyx-mini-0774:phyxminiproblems-problemphyxmini0774-isuniqueroundedanswer}
The assigned autoformalize agent should translate this physics problem into Lean declarations in the covered file, with theorem and lemma proof bodies written as `by sorry`.
\end{theorem}
\begin{proof}
This is an autoformalization task, not a proof task. Produce statements that can later be proved by the physics prover without weakening the physical model.
\end{proof}

% --- Archon declaration topology begin ---
\section{Declaration topology}
\begin{definition}[acceleration Dimension]
  \label{def:physics:phyx-mini-0774:phyxminiproblems-problemphyxmini0774-accelerationdimension}
  \lean{PhyXMiniProblems.ProblemPhyXMini0774.accelerationDimension}
  The dimension L T⁻² of an acceleration magnitude
\end{definition}
\begin{proof}
  This is the declared model definition of acceleration Dimension.
\end{proof}
\begin{definition}[angular Speed Dimension]
  \label{def:physics:phyx-mini-0774:phyxminiproblems-problemphyxmini0774-angularspeeddimension}
  \lean{PhyXMiniProblems.ProblemPhyXMini0774.angularSpeedDimension}
  The dimension T⁻¹ of angular speed; radians are dimensionless
\end{definition}
\begin{proof}
  This is the declared model definition of angular Speed Dimension.
\end{proof}
\begin{definition}[moment Of Inertia Dimension]
  \label{def:physics:phyx-mini-0774:phyxminiproblems-problemphyxmini0774-momentofinertiadimension}
  \lean{PhyXMiniProblems.ProblemPhyXMini0774.momentOfInertiaDimension}
  The dimension M L² of a scalar moment of inertia about an axis
\end{definition}
\begin{proof}
  This is the declared model definition of moment Of Inertia Dimension.
\end{proof}
\begin{definition}[Mass Quantity]
  \label{def:physics:phyx-mini-0774:phyxminiproblems-problemphyxmini0774-massquantity}
  \lean{PhyXMiniProblems.ProblemPhyXMini0774.MassQuantity}
  A nonnegative, unit-independent block mass
\end{definition}
\begin{proof}
  This is the declared model definition of Mass Quantity.
\end{proof}
\begin{definition}[Length Quantity]
  \label{def:physics:phyx-mini-0774:phyxminiproblems-problemphyxmini0774-lengthquantity}
  \lean{PhyXMiniProblems.ProblemPhyXMini0774.LengthQuantity}
  A nonnegative, unit-independent length
\end{definition}
\begin{proof}
  This is the declared model definition of Length Quantity.
\end{proof}
\begin{definition}[Speed Quantity]
  \label{def:physics:phyx-mini-0774:phyxminiproblems-problemphyxmini0774-speedquantity}
  \lean{PhyXMiniProblems.ProblemPhyXMini0774.SpeedQuantity}
  A nonnegative, unit-independent speed magnitude
\end{definition}
\begin{proof}
  This is the declared model definition of Speed Quantity.
\end{proof}
\begin{definition}[Acceleration Quantity]
  \label{def:physics:phyx-mini-0774:phyxminiproblems-problemphyxmini0774-accelerationquantity}
  \lean{PhyXMiniProblems.ProblemPhyXMini0774.AccelerationQuantity}
  \uses{def:physics:phyx-mini-0774:phyxminiproblems-problemphyxmini0774-accelerationdimension}
  A nonnegative, unit-independent acceleration magnitude
\end{definition}
\begin{proof}
  This is the declared model definition of Acceleration Quantity.
\end{proof}
\begin{definition}[Angular Speed Quantity]
  \label{def:physics:phyx-mini-0774:phyxminiproblems-problemphyxmini0774-angularspeedquantity}
  \lean{PhyXMiniProblems.ProblemPhyXMini0774.AngularSpeedQuantity}
  \uses{def:physics:phyx-mini-0774:phyxminiproblems-problemphyxmini0774-angularspeeddimension}
  A nonnegative angular-speed magnitude
\end{definition}
\begin{proof}
  This is the declared model definition of Angular Speed Quantity.
\end{proof}
\begin{definition}[Moment Of Inertia Quantity]
  \label{def:physics:phyx-mini-0774:phyxminiproblems-problemphyxmini0774-momentofinertiaquantity}
  \lean{PhyXMiniProblems.ProblemPhyXMini0774.MomentOfInertiaQuantity}
  \uses{def:physics:phyx-mini-0774:phyxminiproblems-problemphyxmini0774-momentofinertiadimension}
  A nonnegative axial moment of inertia
\end{definition}
\begin{proof}
  This is the declared model definition of Moment Of Inertia Quantity.
\end{proof}
\begin{definition}[Energy Quantity]
  \label{def:physics:phyx-mini-0774:phyxminiproblems-problemphyxmini0774-energyquantity}
  \lean{PhyXMiniProblems.ProblemPhyXMini0774.EnergyQuantity}
  A unit-independent physical energy
\end{definition}
\begin{proof}
  This is the declared model definition of Energy Quantity.
\end{proof}
\begin{definition}[mass In Kilograms]
  \label{def:physics:phyx-mini-0774:phyxminiproblems-problemphyxmini0774-massinkilograms}
  \lean{PhyXMiniProblems.ProblemPhyXMini0774.massInKilograms}
  \uses{def:physics:phyx-mini-0774:phyxminiproblems-problemphyxmini0774-massquantity}
  Read a physical mass in coherent-SI kilograms
\end{definition}
\begin{proof}
  This is the declared model definition of mass In Kilograms.
\end{proof}
\begin{definition}[length In Meters]
  \label{def:physics:phyx-mini-0774:phyxminiproblems-problemphyxmini0774-lengthinmeters}
  \lean{PhyXMiniProblems.ProblemPhyXMini0774.lengthInMeters}
  \uses{def:physics:phyx-mini-0774:phyxminiproblems-problemphyxmini0774-lengthquantity}
  Read a physical length in coherent-SI metres
\end{definition}
\begin{proof}
  This is the declared model definition of length In Meters.
\end{proof}
\begin{definition}[speed In Meters Per Second]
  \label{def:physics:phyx-mini-0774:phyxminiproblems-problemphyxmini0774-speedinmeterspersecond}
  \lean{PhyXMiniProblems.ProblemPhyXMini0774.speedInMetersPerSecond}
  \uses{def:physics:phyx-mini-0774:phyxminiproblems-problemphyxmini0774-speedquantity}
  Read a physical speed in coherent-SI metres per second
\end{definition}
\begin{proof}
  This is the declared model definition of speed In Meters Per Second.
\end{proof}
\begin{definition}[acceleration In Meters Per Second Squared]
  \label{def:physics:phyx-mini-0774:phyxminiproblems-problemphyxmini0774-accelerationinmeterspersecondsquared}
  \lean{PhyXMiniProblems.ProblemPhyXMini0774.accelerationInMetersPerSecondSquared}
  \uses{def:physics:phyx-mini-0774:phyxminiproblems-problemphyxmini0774-accelerationquantity}
  Read an acceleration in coherent-SI metres per second squared
\end{definition}
\begin{proof}
  This is the declared model definition of acceleration In Meters Per Second Squared.
\end{proof}
\begin{definition}[angular Speed In Radians Per Second]
  \label{def:physics:phyx-mini-0774:phyxminiproblems-problemphyxmini0774-angularspeedinradianspersecond}
  \lean{PhyXMiniProblems.ProblemPhyXMini0774.angularSpeedInRadiansPerSecond}
  \uses{def:physics:phyx-mini-0774:phyxminiproblems-problemphyxmini0774-angularspeedquantity}
  Read angular speed in radians per coherent-SI second
\end{definition}
\begin{proof}
  This is the declared model definition of angular Speed In Radians Per Second.
\end{proof}
\begin{definition}[moment Of Inertia In Kilogram Meters Squared]
  \label{def:physics:phyx-mini-0774:phyxminiproblems-problemphyxmini0774-momentofinertiainkilogrammeterssquared}
  \lean{PhyXMiniProblems.ProblemPhyXMini0774.momentOfInertiaInKilogramMetersSquared}
  \uses{def:physics:phyx-mini-0774:phyxminiproblems-problemphyxmini0774-momentofinertiaquantity}
  Read an axial moment of inertia in coherent-SI kg m²
\end{definition}
\begin{proof}
  This is the declared model definition of moment Of Inertia In Kilogram Meters Squared.
\end{proof}
\begin{definition}[energy In Joules]
  \label{def:physics:phyx-mini-0774:phyxminiproblems-problemphyxmini0774-energyinjoules}
  \lean{PhyXMiniProblems.ProblemPhyXMini0774.energyInJoules}
  \uses{def:physics:phyx-mini-0774:phyxminiproblems-problemphyxmini0774-energyquantity}
  Read a physical energy in joules, calibrated by Physlib's joule
\end{definition}
\begin{proof}
  This is the declared model definition of energy In Joules.
\end{proof}
\begin{definition}[Block Label]
  \label{def:physics:phyx-mini-0774:phyxminiproblems-problemphyxmini0774-blocklabel}
  \lean{PhyXMiniProblems.ProblemPhyXMini0774.BlockLabel}
  The two blocks, named by their positions and printed mass labels
\end{definition}
\begin{proof}
  This is the declared model definition of Block Label.
\end{proof}
\begin{definition}[Figure Side]
  \label{def:physics:phyx-mini-0774:phyxminiproblems-problemphyxmini0774-figureside}
  \lean{PhyXMiniProblems.ProblemPhyXMini0774.FigureSide}
  Left/right placement in the supplied diagram
\end{definition}
\begin{proof}
  This is the declared model definition of Figure Side.
\end{proof}
\begin{definition}[Initial Placement]
  \label{def:physics:phyx-mini-0774:phyxminiproblems-problemphyxmini0774-initialplacement}
  \lean{PhyXMiniProblems.ProblemPhyXMini0774.InitialPlacement}
  Initial support relation visible in the supplied diagram
\end{definition}
\begin{proof}
  This is the declared model definition of Initial Placement.
\end{proof}
\begin{definition}[Vertical Direction]
  \label{def:physics:phyx-mini-0774:phyxminiproblems-problemphyxmini0774-verticaldirection}
  \lean{PhyXMiniProblems.ProblemPhyXMini0774.VerticalDirection}
  Vertical directions indicated by the intended motion
\end{definition}
\begin{proof}
  This is the declared model definition of Vertical Direction.
\end{proof}
\begin{definition}[Figure Object]
  \label{def:physics:phyx-mini-0774:phyxminiproblems-problemphyxmini0774-figureobject}
  \lean{PhyXMiniProblems.ProblemPhyXMini0774.FigureObject}
  Physical and graphical objects visible in image 774.png
\end{definition}
\begin{proof}
  This is the declared model definition of Figure Object.
\end{proof}
\begin{definition}[Pulley Mounting]
  \label{def:physics:phyx-mini-0774:phyxminiproblems-problemphyxmini0774-pulleymounting}
  \lean{PhyXMiniProblems.ProblemPhyXMini0774.PulleyMounting}
  How the pulley axle is mounted
\end{definition}
\begin{proof}
  This is the declared model definition of Pulley Mounting.
\end{proof}
\begin{definition}[Rope Model]
  \label{def:physics:phyx-mini-0774:phyxminiproblems-problemphyxmini0774-ropemodel}
  \lean{PhyXMiniProblems.ProblemPhyXMini0774.RopeModel}
  Idealization of the connecting rope used by the textbook model
\end{definition}
\begin{proof}
  This is the declared model definition of Rope Model.
\end{proof}
\begin{definition}[Rope Pulley Contact]
  \label{def:physics:phyx-mini-0774:phyxminiproblems-problemphyxmini0774-ropepulleycontact}
  \lean{PhyXMiniProblems.ProblemPhyXMini0774.RopePulleyContact}
  Contact condition between the rope and the pulley rim
\end{definition}
\begin{proof}
  This is the declared model definition of Rope Pulley Contact.
\end{proof}
\begin{definition}[Axle Model]
  \label{def:physics:phyx-mini-0774:phyxminiproblems-problemphyxmini0774-axlemodel}
  \lean{PhyXMiniProblems.ProblemPhyXMini0774.AxleModel}
  Dissipation model at the fixed axle
\end{definition}
\begin{proof}
  This is the declared model definition of Axle Model.
\end{proof}
\begin{definition}[Release State]
  \label{def:physics:phyx-mini-0774:phyxminiproblems-problemphyxmini0774-releasestate}
  \lean{PhyXMiniProblems.ProblemPhyXMini0774.ReleaseState}
  Initial state needed by the recorded energy calculation
\end{definition}
\begin{proof}
  This is the declared model definition of Release State.
\end{proof}
\begin{definition}[Measurement Event]
  \label{def:physics:phyx-mini-0774:phyxminiproblems-problemphyxmini0774-measurementevent}
  \lean{PhyXMiniProblems.ProblemPhyXMini0774.MeasurementEvent}
  The event at which the requested speed is measured
\end{definition}
\begin{proof}
  This is the declared model definition of Measurement Event.
\end{proof}
\begin{definition}[Massive Pulley Figure]
  \label{def:physics:phyx-mini-0774:phyxminiproblems-problemphyxmini0774-massivepulleyfigure}
  \lean{PhyXMiniProblems.ProblemPhyXMini0774.MassivePulleyFigure}
  \uses{def:physics:phyx-mini-0774:phyxminiproblems-problemphyxmini0774-blocklabel, def:physics:phyx-mini-0774:phyxminiproblems-problemphyxmini0774-figureside, def:physics:phyx-mini-0774:phyxminiproblems-problemphyxmini0774-initialplacement, def:physics:phyx-mini-0774:phyxminiproblems-problemphyxmini0774-verticaldirection, def:physics:phyx-mini-0774:phyxminiproblems-problemphyxmini0774-figureobject}
  Literal qualitative and numerical information transcribed from the primary bitmap. The ordering data describe the drawing; they are not physical scalar coordinates
\end{definition}
\begin{proof}
  This is the declared model definition of Massive Pulley Figure.
\end{proof}
\begin{definition}[Massive Pulley Atwood Setup]
  \label{def:physics:phyx-mini-0774:phyxminiproblems-problemphyxmini0774-massivepulleyatwoodsetup}
  \lean{PhyXMiniProblems.ProblemPhyXMini0774.MassivePulleyAtwoodSetup}
  \uses{def:physics:phyx-mini-0774:phyxminiproblems-problemphyxmini0774-massquantity, def:physics:phyx-mini-0774:phyxminiproblems-problemphyxmini0774-lengthquantity, def:physics:phyx-mini-0774:phyxminiproblems-problemphyxmini0774-speedquantity, def:physics:phyx-mini-0774:phyxminiproblems-problemphyxmini0774-accelerationquantity, def:physics:phyx-mini-0774:phyxminiproblems-problemphyxmini0774-angularspeedquantity, def:physics:phyx-mini-0774:phyxminiproblems-problemphyxmini0774-momentofinertiaquantity, def:physics:phyx-mini-0774:phyxminiproblems-problemphyxmini0774-energyquantity, def:physics:phyx-mini-0774:phyxminiproblems-problemphyxmini0774-blocklabel, def:physics:phyx-mini-0774:phyxminiproblems-problemphyxmini0774-verticaldirection, def:physics:phyx-mini-0774:phyxminiproblems-problemphyxmini0774-pulleymounting, def:physics:phyx-mini-0774:phyxminiproblems-problemphyxmini0774-ropemodel, def:physics:phyx-mini-0774:phyxminiproblems-problemphyxmini0774-ropepulleycontact, def:physics:phyx-mini-0774:phyxminiproblems-problemphyxmini0774-axlemodel, def:physics:phyx-mini-0774:phyxminiproblems-problemphyxmini0774-releasestate, def:physics:phyx-mini-0774:phyxminiproblems-problemphyxmini0774-measurementevent, def:physics:phyx-mini-0774:phyxminiproblems-problemphyxmini0774-massivepulleyfigure}
  Independent physical quantities and qualitative model choices. The scalar impact speed and angular speed are fields rather than definitions made from the recorded answer or from the target closed form
\end{definition}
\begin{proof}
  This is the declared model definition of Massive Pulley Atwood Setup.
\end{proof}
\begin{definition}[Matches Primary Figure]
  \label{def:physics:phyx-mini-0774:phyxminiproblems-problemphyxmini0774-matchesprimaryfigure}
  \lean{PhyXMiniProblems.ProblemPhyXMini0774.MatchesPrimaryFigure}
  \uses{def:physics:phyx-mini-0774:phyxminiproblems-problemphyxmini0774-massinkilograms, def:physics:phyx-mini-0774:phyxminiproblems-problemphyxmini0774-lengthinmeters, def:physics:phyx-mini-0774:phyxminiproblems-problemphyxmini0774-massivepulleyatwoodsetup}
  Geometry and labels read directly from the primary image
\end{definition}
\begin{proof}
  This is the declared model definition of Matches Primary Figure.
\end{proof}
\begin{definition}[Matches Problem Statement]
  \label{def:physics:phyx-mini-0774:phyxminiproblems-problemphyxmini0774-matchesproblemstatement}
  \lean{PhyXMiniProblems.ProblemPhyXMini0774.MatchesProblemStatement}
  \uses{def:physics:phyx-mini-0774:phyxminiproblems-problemphyxmini0774-lengthinmeters, def:physics:phyx-mini-0774:phyxminiproblems-problemphyxmini0774-momentofinertiainkilogrammeterssquared, def:physics:phyx-mini-0774:phyxminiproblems-problemphyxmini0774-massivepulleyatwoodsetup}
  The numerical pulley data and no-slip statement from the prose. No impact speed or answer-choice value occurs in these premises
\end{definition}
\begin{proof}
  This is the declared model definition of Matches Problem Statement.
\end{proof}
\begin{definition}[Uses Ideal Textbook Release Model]
  \label{def:physics:phyx-mini-0774:phyxminiproblems-problemphyxmini0774-usesidealtextbookreleasemodel}
  \lean{PhyXMiniProblems.ProblemPhyXMini0774.UsesIdealTextbookReleaseModel}
  \uses{def:physics:phyx-mini-0774:phyxminiproblems-problemphyxmini0774-massivepulleyatwoodsetup}
  Additional idealizations needed by the textbook answer but not stated explicitly in the short prompt: release from rest, a massless inextensible rope, and no dissipative axle friction
\end{definition}
\begin{proof}
  This is the declared model definition of Uses Ideal Textbook Release Model.
\end{proof}
\begin{definition}[Uses Standard Near Earth Gravity]
  \label{def:physics:phyx-mini-0774:phyxminiproblems-problemphyxmini0774-usesstandardnearearthgravity}
  \lean{PhyXMiniProblems.ProblemPhyXMini0774.UsesStandardNearEarthGravity}
  \uses{def:physics:phyx-mini-0774:phyxminiproblems-problemphyxmini0774-accelerationinmeterspersecondsquared, def:physics:phyx-mini-0774:phyxminiproblems-problemphyxmini0774-massivepulleyatwoodsetup}
  Standard near-Earth gravitational calibration used by the recorded answer
\end{definition}
\begin{proof}
  This is the declared model definition of Uses Standard Near Earth Gravity.
\end{proof}
\begin{definition}[Has Physical Parameters]
  \label{def:physics:phyx-mini-0774:phyxminiproblems-problemphyxmini0774-hasphysicalparameters}
  \lean{PhyXMiniProblems.ProblemPhyXMini0774.HasPhysicalParameters}
  \uses{def:physics:phyx-mini-0774:phyxminiproblems-problemphyxmini0774-massinkilograms, def:physics:phyx-mini-0774:phyxminiproblems-problemphyxmini0774-lengthinmeters, def:physics:phyx-mini-0774:phyxminiproblems-problemphyxmini0774-speedinmeterspersecond, def:physics:phyx-mini-0774:phyxminiproblems-problemphyxmini0774-accelerationinmeterspersecondsquared, def:physics:phyx-mini-0774:phyxminiproblems-problemphyxmini0774-angularspeedinradianspersecond, def:physics:phyx-mini-0774:phyxminiproblems-problemphyxmini0774-momentofinertiainkilogrammeterssquared, def:physics:phyx-mini-0774:phyxminiproblems-problemphyxmini0774-massivepulleyatwoodsetup}
  Positivity and mass ordering selecting the depicted physical branch
\end{definition}
\begin{proof}
  This is the declared model definition of Has Physical Parameters.
\end{proof}
\begin{definition}[Satisfies Massive Pulley Energy Laws]
  \label{def:physics:phyx-mini-0774:phyxminiproblems-problemphyxmini0774-satisfiesmassivepulleyenergylaws}
  \lean{PhyXMiniProblems.ProblemPhyXMini0774.SatisfiesMassivePulleyEnergyLaws}
  \uses{def:physics:phyx-mini-0774:phyxminiproblems-problemphyxmini0774-massinkilograms, def:physics:phyx-mini-0774:phyxminiproblems-problemphyxmini0774-lengthinmeters, def:physics:phyx-mini-0774:phyxminiproblems-problemphyxmini0774-speedinmeterspersecond, def:physics:phyx-mini-0774:phyxminiproblems-problemphyxmini0774-accelerationinmeterspersecondsquared, def:physics:phyx-mini-0774:phyxminiproblems-problemphyxmini0774-angularspeedinradianspersecond, def:physics:phyx-mini-0774:phyxminiproblems-problemphyxmini0774-momentofinertiainkilogrammeterssquared, def:physics:phyx-mini-0774:phyxminiproblems-problemphyxmini0774-energyinjoules, def:physics:phyx-mini-0774:phyxminiproblems-problemphyxmini0774-massivepulleyatwoodsetup}
  Governing laws for the ideal massive-pulley Atwood model. No slip equates the rope speed to radius times angular speed. The pulley inertia is the tensor entry about its fixed axis, and its angular velocity points along that axis. Pulley rotational energy is Physlib's RigidBody.rotationalKineticEnergy. Conservation of energy equates the net gravitational potential-energy loss to the two blocks' translational kinetic energy plus the pulley's rotational kinetic energy.
\end{definition}
\begin{proof}
  This is the declared model definition of Satisfies Massive Pulley Energy Laws.
\end{proof}
\begin{lemma}[pulley Rotational Energy eq half inertia mul angular Speed sq]
  \label{lem:physics:phyx-mini-0774:phyxminiproblems-problemphyxmini0774-pulleyrotationalenergy-eq-half-inertia-mul-angularspeed-sq}
  \lean{PhyXMiniProblems.ProblemPhyXMini0774.pulleyRotationalEnergy_eq_half_inertia_mul_angularSpeed_sq}
  \uses{def:physics:phyx-mini-0774:phyxminiproblems-problemphyxmini0774-angularspeedinradianspersecond, def:physics:phyx-mini-0774:phyxminiproblems-problemphyxmini0774-momentofinertiainkilogrammeterssquared, def:physics:phyx-mini-0774:phyxminiproblems-problemphyxmini0774-energyinjoules, def:physics:phyx-mini-0774:phyxminiproblems-problemphyxmini0774-massivepulleyatwoodsetup, def:physics:phyx-mini-0774:phyxminiproblems-problemphyxmini0774-satisfiesmassivepulleyenergylaws}
  The tensor contraction reduces to the scalar pulley relation K\_rot = I * omega² / 2 because angular velocity lies along the fixed axis
\end{lemma}
\begin{proof}
  The conclusion follows from the governing relations and figure readouts named in the statement of pulley Rotational Energy eq half inertia mul angular Speed sq.
\end{proof}
\begin{lemma}[impact Speed squared eq 6272 over 667]
  \label{lem:physics:phyx-mini-0774:phyxminiproblems-problemphyxmini0774-impactspeed-squared-eq-6272-over-667}
  \lean{PhyXMiniProblems.ProblemPhyXMini0774.impactSpeed_squared_eq_6272_over_667}
  \uses{def:physics:phyx-mini-0774:phyxminiproblems-problemphyxmini0774-speedinmeterspersecond, def:physics:phyx-mini-0774:phyxminiproblems-problemphyxmini0774-massivepulleyatwoodsetup, def:physics:phyx-mini-0774:phyxminiproblems-problemphyxmini0774-matchesprimaryfigure, def:physics:phyx-mini-0774:phyxminiproblems-problemphyxmini0774-matchesproblemstatement, def:physics:phyx-mini-0774:phyxminiproblems-problemphyxmini0774-usesidealtextbookreleasemodel, def:physics:phyx-mini-0774:phyxminiproblems-problemphyxmini0774-usesstandardnearearthgravity, def:physics:phyx-mini-0774:phyxminiproblems-problemphyxmini0774-hasphysicalparameters, def:physics:phyx-mini-0774:phyxminiproblems-problemphyxmini0774-satisfiesmassivepulleyenergylaws}
  Substitution of the figure, pulley, gravity, no-slip, and energy data gives the exact squared impact speed 6272 / 667 (m/s)²
\end{lemma}
\begin{proof}
  The conclusion follows from the governing relations and figure readouts named in the statement of impact Speed squared eq 6272 over 667.
\end{proof}
\begin{lemma}[impact Speed eq sqrt 6272 over 667]
  \label{lem:physics:phyx-mini-0774:phyxminiproblems-problemphyxmini0774-impactspeed-eq-sqrt-6272-over-667}
  \lean{PhyXMiniProblems.ProblemPhyXMini0774.impactSpeed_eq_sqrt_6272_over_667}
  \uses{def:physics:phyx-mini-0774:phyxminiproblems-problemphyxmini0774-speedinmeterspersecond, def:physics:phyx-mini-0774:phyxminiproblems-problemphyxmini0774-massivepulleyatwoodsetup, def:physics:phyx-mini-0774:phyxminiproblems-problemphyxmini0774-matchesprimaryfigure, def:physics:phyx-mini-0774:phyxminiproblems-problemphyxmini0774-matchesproblemstatement, def:physics:phyx-mini-0774:phyxminiproblems-problemphyxmini0774-usesidealtextbookreleasemodel, def:physics:phyx-mini-0774:phyxminiproblems-problemphyxmini0774-usesstandardnearearthgravity, def:physics:phyx-mini-0774:phyxminiproblems-problemphyxmini0774-hasphysicalparameters, def:physics:phyx-mini-0774:phyxminiproblems-problemphyxmini0774-satisfiesmassivepulleyenergylaws}
  The positive physical branch of the exact impact-speed solution
\end{lemma}
\begin{proof}
  The conclusion follows from the governing relations and figure readouts named in the statement of impact Speed eq sqrt 6272 over 667.
\end{proof}
\begin{definition}[Answer Choice]
  \label{def:physics:phyx-mini-0774:phyxminiproblems-problemphyxmini0774-answerchoice}
  \lean{PhyXMiniProblems.ProblemPhyXMini0774.AnswerChoice}
  Labels of the four speed choices printed in the problem
\end{definition}
\begin{proof}
  This is the declared model definition of Answer Choice.
\end{proof}
\begin{definition}[displayed Speed Meters Per Second]
  \label{def:physics:phyx-mini-0774:phyxminiproblems-problemphyxmini0774-answerchoice-displayedspeedmeterspersecond}
  \lean{PhyXMiniProblems.ProblemPhyXMini0774.AnswerChoice.displayedSpeedMetersPerSecond}
  \uses{def:physics:phyx-mini-0774:phyxminiproblems-problemphyxmini0774-answerchoice}
  Metre-per-second value printed beside each answer label
\end{definition}
\begin{proof}
  This is the declared model definition of displayed Speed Meters Per Second.
\end{proof}
\begin{definition}[recorded Dataset Answer]
  \label{def:physics:phyx-mini-0774:phyxminiproblems-problemphyxmini0774-recordeddatasetanswer}
  \lean{PhyXMiniProblems.ProblemPhyXMini0774.recordedDatasetAnswer}
  \uses{def:physics:phyx-mini-0774:phyxminiproblems-problemphyxmini0774-answerchoice}
  Dataset answer metadata, deliberately not used as a theorem premise
\end{definition}
\begin{proof}
  This is the declared model definition of recorded Dataset Answer.
\end{proof}
\begin{definition}[Rounds To Displayed Speed]
  \label{def:physics:phyx-mini-0774:phyxminiproblems-problemphyxmini0774-roundstodisplayedspeed}
  \lean{PhyXMiniProblems.ProblemPhyXMini0774.RoundsToDisplayedSpeed}
  \uses{def:physics:phyx-mini-0774:phyxminiproblems-problemphyxmini0774-speedinmeterspersecond, def:physics:phyx-mini-0774:phyxminiproblems-problemphyxmini0774-massivepulleyatwoodsetup, def:physics:phyx-mini-0774:phyxminiproblems-problemphyxmini0774-answerchoice}
  A computed speed matches a two-decimal answer when it lies in that displayed hundredth's half-open rounding interval
\end{definition}
\begin{proof}
  This is the declared model definition of Rounds To Displayed Speed.
\end{proof}
\begin{definition}[Is Unique Rounded Answer]
  \label{def:physics:phyx-mini-0774:phyxminiproblems-problemphyxmini0774-isuniqueroundedanswer}
  \lean{PhyXMiniProblems.ProblemPhyXMini0774.IsUniqueRoundedAnswer}
  \uses{def:physics:phyx-mini-0774:phyxminiproblems-problemphyxmini0774-massivepulleyatwoodsetup, def:physics:phyx-mini-0774:phyxminiproblems-problemphyxmini0774-answerchoice, def:physics:phyx-mini-0774:phyxminiproblems-problemphyxmini0774-roundstodisplayedspeed}
  A choice is the unique displayed speed to which the impact speed rounds
\end{definition}
\begin{proof}
  This is the declared model definition of Is Unique Rounded Answer.
\end{proof}
% --- Archon declaration topology end ---
% --- Archon physics formalization source end ---
